\documentclass[12pt]{article}
\usepackage[T1]{fontenc}
\usepackage[utf8]{inputenc}
\usepackage{lmodern}
\usepackage{textcomp}
\usepackage{enumitem}
\usepackage{geometry}
\geometry{margin=1in}
\begin{document}
\begin{center}
Answer Key - Biology - Graduate Level Assessment 18 \
\small (Answers are ordered exactly as in the question paper.)
\end{center}

\section*{Multiple Choice}
\begin{enumerate}[label=\arabic*.]
\item \textbf{Answer:} A
\\ \textit{Options:} A) 50\%; B) 5\%; C) 60\%; D) 25\%; E) 90\%
\\ \textit{Note:} Correct option: A - 50\%
\item \textbf{Answer:} C
\\ \textit{Options:} A) The main source of variation in humans is genetic engineering.; B) The main source of variation in humans is due to dietary factors.; C) Shuffling of alleles is responsible for the greatest variation in humans.; D) The main source of variation in humans is due to exposure to radiation.; E) The main source of variation in humans is genetic drift.
\\ \textit{Note:} Correct option: C - Shuffling of alleles is responsible for the greatest variation in humans.
\item \textbf{Answer:} A
\\ \textit{Options:} A) The number of synapses crossed; B) The speed of the action potential; C) The amplitude of the action potential; D) The number of neurons activated; E) The refractory period of the neuron
\\ \textit{Note:} Correct option: A - The number of synapses crossed
\item \textbf{Answer:} E
\\ \textit{Options:} A) - .09176; B) .21185; C) .37748; D) .55962; E) .65361
\\ \textit{Note:} Correct option: E - .65361
\item \textbf{Answer:} C
\\ \textit{Options:} A) parenchyma; B) phloem; C) endoderm; D) pericycle; E) cortex
\\ \textit{Note:} Correct option: C - endoderm
\end{enumerate}

\section*{True / False}
\begin{enumerate}[label=\arabic*.]
\setcounter{enumi}{5}
\item \textbf{Answer:} False
\\ \textit{Options:} A) True; B) False
\\ \textit{Note:} Clinical variables, either isolated or as components of a model, could not identify all children with pathologic radiographs.
\item \textbf{Answer:} True
\\ \textit{Options:} A) True; B) False
\\ \textit{Note:} Both medial and lateral PTS were significantly steeper in failures of ACLR than the control group. Medial or lateral PTS \ensuremath{\geq} 5^\ensuremath{\circ} was a new risk factor of ACLR failure.
\item \textbf{Answer:} True
\\ \textit{Options:} A) True; B) False
\\ \textit{Note:} These findings show that phagocytic NADPH oxidase activity is increased in obesity and is related to preclinical atherosclerosis in this condition. We also suggest that hyperleptinemia may contribute to phagocytic NADPH oxidase overactivity in obesity.
\item \textbf{Answer:} True
\\ \textit{Options:} A) True; B) False
\\ \textit{Note:} Speed discrimination, per se, is not impaired in schizophrenia patients. The observed abnormality appears to be a consequence of impairment in generating or integrating the feedback information from eye movements. This study introduces a novel approach to motion perception studies and highlights the importance of concurrently measuring eye movements to understand interactions between these two systems; the results argue for a conceptual revision regarding motion perception abnormality in schizophrenia.
\item \textbf{Answer:} False
\\ \textit{Options:} A) True; B) False
\\ \textit{Note:} Patients with WD may possibly undergo cardiac surgery without a markedly enhanced risk for bleeding complications despite a more than usual transfusion requirement and significantly lower platelet counts perioperatively.
\end{enumerate}

\section*{Long Form Response}
\begin{enumerate}[label=\arabic*.]
\setcounter{enumi}{10}
\item \textbf{Answer:} 60 possible genotypes. Explain why this is the correct answer and provide supporting evidence. Reference key facts or concepts that support this choice.
\\ \textit{Note:} Long-form derived from MCQ; award credit for stating the correct idea and giving supporting reasoning.
\item \textbf{Answer:} Cross-pollination is promoted only through manual intervention by humans, without any natural methods.. Explain why this is the correct answer and provide supporting evidence. Reference key facts or concepts that support this choice.
\\ \textit{Note:} Long-form derived from MCQ; award credit for stating the correct idea and giving supporting reasoning.
\end{enumerate}

\vfill
\noindent\textit{End of Answer Key}
\end{document}